\documentclass[11pt,reqno]{amsart}

\usepackage[T1]{fontenc}
\usepackage[utf8]{inputenc}
\usepackage{newtxtext,newtxmath}
\usepackage[a4paper,margin=30mm]{geometry}
\usepackage{microtype}
\usepackage{amsmath,amsthm,mathtools}
\usepackage{bm,mathrsfs}
\usepackage{booktabs,longtable,tabularx,array}
\usepackage{enumitem}
\usepackage[dvipsnames]{xcolor}
\usepackage{xurl}
\usepackage[hidelinks]{hyperref}
\usepackage[nameinlink,capitalize,noabbrev]{cleveref}
\usepackage[backend=biber,style=alphabetic,maxbibnames=99,giveninits=true]{biblatex}
\addbibresource{twin_gate1b_r11_audited_20260905.bib}

\allowdisplaybreaks
\setlength{\emergencystretch}{3em}
\setcounter{tocdepth}{2}
\numberwithin{equation}{section}

\newtheorem{theorem}{Theorem}[section]
\newtheorem{proposition}[theorem]{Proposition}
\newtheorem{lemma}[theorem]{Lemma}
\newtheorem{corollary}[theorem]{Corollary}
\newtheorem{conjecture}[theorem]{Candidate theorem}
\newtheorem{problem}[theorem]{Open problem}
\theoremstyle{definition}
\newtheorem{definition}[theorem]{Definition}
\newtheorem{hypothesis}[theorem]{Frozen-bank hypothesis}
\newtheorem{status}[theorem]{Status statement}
\theoremstyle{remark}
\newtheorem{remark}[theorem]{Remark}
\newtheorem{warning}[theorem]{Firewall}

\newcommand{\e}{\mathrm e}
\newcommand{\ee}[1]{\e\!\left(#1\right)}
\newcommand{\eord}[2]{\e_{#1}\!\left(#2\right)}
\newcommand{\1}{\mathbf 1}
\newcommand{\C}{\mathbb C}
\newcommand{\R}{\mathbb R}
\newcommand{\Z}{\mathbb Z}
\newcommand{\N}{\mathbb N}
\newcommand{\T}{\mathbb T}
\newcommand{\LL}{\mathcal L}
\newcommand{\RR}{\mathcal R}
\newcommand{\BB}{\mathcal B}
\newcommand{\cA}{\mathcal A}
\newcommand{\cC}{\mathcal C}
\newcommand{\cE}{\mathcal E}
\newcommand{\cI}{\mathcal I}
\newcommand{\cP}{\mathcal P}
\newcommand{\cO}{\mathcal O}
\newcommand{\wt}{\widetilde}
\newcommand{\wh}{\widehat}
\newcommand{\norm}[1]{\left\lVert #1\right\rVert}
\newcommand{\abs}[1]{\left|#1\right|}
\newcommand{\ip}[2]{\left\langle #1,#2\right\rangle}
\newcommand{\Om}{\Omega}
\newcommand{\eps}{\varepsilon}
\newcommand{\ph}{\varphi}
\newcommand{\Rhi}{\RR_{\mathrm{hi}}}
\newcommand{\Roneone}{\RR_{11}}
\newcommand{\soft}{X^{o(1)}}
\newcommand{\Pmin}{P_*}
\newcommand{\epsstar}{\eps_*}
\newcommand{\epssafe}{\eps_{\mathrm{safe}}}

\newcolumntype{P}[1]{>{\raggedright\arraybackslash}p{#1}}
\newcommand{\classification}[2]{%
  \par\smallskip\noindent
  {\footnotesize\textbf{Novelty classification:} \textsc{#1}.\quad
   \textbf{Evidence level:} #2.}\par\smallskip}
\newcommand{\standardresult}{STANDARD RESULT}
\newcommand{\newderivation}{NEW DERIVATION}
\newcommand{\newapplication}{NEW SOURCE-SPECIFIC APPLICATION}
\newcommand{\newreduction}{NEW EXACT REDUCTION}
\newcommand{\candidatenew}{CANDIDATE NEW THEOREM}
\newcommand{\openstatus}{OPEN}

\newenvironment{statusbox}{%
  \par\medskip\noindent\begin{minipage}{\textwidth}\hrule\medskip\small
}{%
  \par\medskip\hrule\end{minipage}\par\medskip
}

\title[The Gate 1B $R_{11}$ frontier]{An Audited Rough-Conductor Reduction for the Gate 1B Eleven-Prime Residual}
\author{Nam Chong Lau}
\address{Independent research programme}
\email{}
\date{5 September 2026}
\subjclass[2020]{11N35, 11N36, 11L40, 11N75, 68V20}
\keywords{twin-prime sieve, Type II sums, multiplicative large sieve, rough moduli, M\"obius function, finite Fourier analysis, formal verification}

\begin{document}
\begin{abstract}
We give a source-explicit and audit-separated account of the current Gate~1B
$R_{11}$ programme.  The purpose is not to announce an $R_{11}$ estimate, a
closure of Gate~1B, an $H^*$ theorem, or a proof of the twin-prime conjecture.
Rather, we identify exactly which parts of the reduction are algebraically
proved, which analytic estimates have survived independent hostile audit, which
steps remain inherited hypotheses, and which single physical functional is
presently open.

At the exact level, an eleven-coordinate ordered exponent geometry gives the
uniform cutoff $M\le 5/11$ and complementary mass at least $6/11$.  The
$V=2$ Vaughan chart agrees literally with the long-M\"obius chart on odd
arguments; the CARD5 coefficient is $\binom{10}{5}=252$; the internal
$4\mid4\mid2$ allocation count $3150$ is completely consumed by its
normalisation; and the matched equation
\[
        AB-kd=-2
\]
admits the unique affine parametrisation
\[
        d_t=d_0+At,\qquad B_t=B_0+kt.
\]
The associated two-copy determinant is
\[
        d_tB_{t+h}-d_{t+h}B_t=2h.
\]
These finite identities, together with an abstract Hilbert-valued reciprocal
large-sieve theorem, are machine checked in Lean.  The formal development does
not construct the physical Hilbert caller and does not prove that the canonical
source equals the full historical $R_{11}$ source.

The new analytic input is a rough-modulus variance estimate.  If
$q\in[R,2R)$ is squarefree, has at most $j$ prime factors, and has least prime
factor at least $P$, then for a single common $1$-bounded sequence supported on
$(D,2D]$,
\[
 \sum_q\sum_{a\bmod q}^{*}\left|
 \sum_{\substack{D<n\le y\\n\equiv a\, (q)}}a_n
 -\frac1{\phi(q)}\sum_{\substack{D<n\le y\\(n,q)=1}}a_n
 \right|^2
 \ll_j DR+\frac{D^2}{P}.
\]
We prove this from the primitive multiplicative large sieve while retaining the
full imprimitive coprimality correction.  The theorem transfers to the literal
band-limited six-prime target and closes the former balanced-low Sector~A at
arbitrary logarithmic precision.  It also closes every fixed-polylogarithmic
frequency band throughout the larger region $K\ge R/\sqrt{P_*}$, and further
where an explicit cellwise conductor deficit allows.

After the independently audited corrections, the principal coprimality cost is
$O_j(2^j(D/P+1))$, the conservative logarithmic ledger is
$\LL^{114+3b/2}$ rather than $\LL^{113+3b/2}$, and the frequency cutoff is
formed using the minimum cyclic conductor in each cell.  The old three-sector
FSHC formulation is therefore superseded by the exact first-moment identity
\[
        \Roneone(X)=\Rhi(X)+O_S(X\LL^{-S}),
\]
where $\Rhi$ contains only the unremoved target frequencies and retains the
literal M\"obius coordinate, all six-prime labels, and the complete physical
normalisation.  No bound of the required strength is presently known for
$\Rhi$.  Thus $R_{11}$ remains open.
\end{abstract}

\maketitle

\begin{statusbox}
\textbf{Authoritative status on 5 September 2026.}
The former Sector~A is closed.  The former three-sector FSHC formulation is
superseded.  The physical delta-plus-single-Poisson Hilbert--HMRD caller is
retired as a nonclosing architecture, while the finite abstract Hilbert--HMRD
theorem remains valid.  The current physical residual is $\Rhi$.  The nodes
$R_{11}$, Global Gate~1B, $H^*$, and the twin-prime conjecture are all open.
\end{statusbox}

\tableofcontents

\section{Introduction}

The analytic difficulty addressed here is not the manufacture of another
arbitrary-coefficient Type~II estimate.  The relevant object is a highly
structured, signed, source-generated functional arising after an eleven-prime
sieve source has been partitioned, reassembled, and placed on a determinant
line.  Every useful estimate must therefore preserve four kinds of information
simultaneously: the five-prime modulus, the six-prime target, the affine relation
between the M\"obius and target variables, and the finite normalisation inherited
from the source compiler.

This manuscript has two aims.  The first is expository: to replace a sequence
of research ledgers by a reviewable theorem--lemma--proof narrative.  The second
is epistemic: to keep distinct the three evidentiary levels that are often
silently conflated in long computational research programmes.

\begin{enumerate}[label=\textup{(E\arabic*)},leftmargin=3em]
\item \textbf{Kernel-checked exact mathematics.}  These are finite identities,
      inequalities, and conditional interfaces proved in Lean~4 and mathlib
      \cite{Lean4CADE2021,Mathlib2020}.
\item \textbf{Independently audited analytic mathematics.}  These are paper
      arguments replayed line by line by an independent hostile audit.  They
      are not yet formalised in Lean.
\item \textbf{Frozen source hypotheses.}  These are prior source-compilation,
      coverage, and local-sector claims used as hypotheses in the newest run.
      The present manuscript records but does not silently re-prove them.
\end{enumerate}

This separation matters.  In particular, the formal bank proves a canonical
CARD5 source and proves many exact properties of it, but deliberately leaves
\emph{CanonicalCoversFullR11} uninhabited.  Conversely, the analytic
rough-conductor argument is conducted on the supplied full physical source,
but its identification with historical $R_{11}$ is inherited rather than
kernel checked.  The results below are stated with precisely those dependencies.

The broad downstream motivation is the prime-producing sieve framework of Ford
and Maynard \cite{FordMaynard2024}.  Nothing in this paper verifies the entire
downstream compiler.  Consequently, even a future proof of the displayed
high-pass estimate would first close the present $R_{11}$ node; it would not by
itself constitute a proof of twin primes.

\subsection{What has changed relative to the previous frontier}

The previous frequency-spliced MHHC analysis removed the diagonal, a
near-diagonal range, large determinant-gcd strata, exceptional five-prime
moduli, and explicit pretentious character modes.  Its balanced-low target
estimate supplied only a fixed logarithmic gain.  The rough-modulus variance
theorem developed below removes the dependence on that weak balanced-low
estimate and proves arbitrary-logarithmic first-moment control in a larger
region.  The correct end product is no longer a three-piece absolute-value
problem.  It is one signed high-pass functional $\Rhi$ obtained by a linear
projection before any new two-copy expansion.

\subsection{Novelty discipline}

Every main result is labelled using one of the categories requested for this
research bank.  ``New derivation'' means that the displayed argument was newly
derived in this programme; it does not assert literature priority.  ``New
source-specific application'' means that a standard or newly derived tool is
inserted into the literal physical source.  ``Candidate new theorem'' is a
research target, not an established result.  The labels are intentionally
conservative.

\section{Main results and the audited frontier}\label{sec:main}

We begin by isolating the inherited assumptions under which the newest analytic
claims are made.

\begin{hypothesis}[Frozen physical source bank]\label{hyp:frozen}
For each dyadic cell $C$, the supplied $R_{11}$ compiler produces variables
\[
 A\asymp R_C,\qquad k\asymp K_C,\qquad d\asymp D_C,
 \qquad B\asymp \BB_C,
\]
with
\[
 D_CK_C\asymp X,\qquad R_C T_C\asymp D_C,
 \qquad K_CT_C\asymp\BB_C,
\]
and with the matched relation $AB-kd=-2$.  The modulus $A$ is squarefree,
has exactly five physical prime factors, and satisfies
\[
       P^-(A)\ge P_C\ge P_*\gg X^{2\epsstar}
\]
for a fixed $\epsstar>0$.  The target is the literal labelled six-prime
coefficient described in \cref{sec:physical}.  The low-divisor compiler,
previously owned local sectors, and the finite source reassembly contribute
